\documentclass[aspectratio=169]{beamer}
\usetheme{Madrid}
\usecolortheme{whale}

% --- Configuration des packages ---
\usepackage[utf8]{inputenc}
\usepackage[T1]{fontenc}
\usepackage[french]{babel}
\usepackage{graphicx}
\usepackage{amsmath}
\usepackage{booktabs}
\usepackage{hyperref}
\usepackage{tikz}

% --- Personnalisation des couleurs (Bleu/Cyan Moderne) ---
\definecolor{maincolor}{RGB}{0, 102, 204}
\definecolor{accentcolor}{RGB}{0, 153, 153}
\setbeamercolor{structure}{fg=maincolor}
\setbeamercolor{palette primary}{bg=maincolor, fg=white}
\setbeamercolor{palette secondary}{bg=accentcolor, fg=white}
\setbeamercolor{palette tertiary}{bg=darkgray, fg=white}

% --- Métadonnées ---
\title[Archi-Agent VR]{Archi-Agent VR}
\subtitle{Assistant Architectural Multimodal VR \& PC par la Voix}
\author[Présentation]{Présentation du Projet}
\institute[Université]{Recherche en Interaction Multimodale}
\date{Juin 2026}

\begin{document}

% ==========================================
% SLIDE 1 : Page de titre
% ==========================================
\begin{frame}
  \titlepage
  
  % %%%%%%%%%%%%%%%%%%%%%%%%%%%%%%%%%%%%%%%%%%
  % % === À DIRE À L'ORAL ===
  % % "Bonjour à tous. Je vais vous présenter aujourd'hui le projet Archi-Agent VR. 
  % % C'est un assistant d'architecture et de design d'intérieur immersif, utilisable en réalité virtuelle 
  % % et sur PC. Notre objectif principal a été de concevoir une interface où l'utilisateur pilote la 
  % % création de son espace 3D uniquement par la voix, en langage naturel, avec un retour fluide en 
  % % moins de 5 secondes."
  % %%%%%%%%%%%%%%%%%%%%%%%%%%%%%%%%%%%%%%%%%%
\end{frame}


% ==========================================
% SLIDE 2 : Problématique et Objectifs
% ==========================================
\begin{frame}{📖 Contexte et Objectifs}
  \begin{columns}
    \begin{column}{0.5\textwidth}
      \textbf{La Problématique :}
      \begin{itemize}
        \item Complexité des outils CAO traditionnels.
        \item Fatigue en VR avec la navigation aux joysticks.
        \item Interfaces 2D intrusives en immersion 3D.
      \end{itemize}
    \end{column}
    \begin{column}{0.5\textwidth}
      \textbf{Nos Objectifs :}
      \begin{itemize}
        \item \textbf{Zéro manette pour la navigation} : La voix gère les déplacements et la construction.
        \item \textbf{Orchestration LLM} : Llama 3.3 \& Gemini pilotent le monde 3D via des appels d'outils.
        \item \textbf{Haute disponibilité} : Cascades d'APIs gratuites pour une résilience totale.
      \end{itemize}
    \end{column}
  \end{columns}

  % %%%%%%%%%%%%%%%%%%%%%%%%%%%%%%%%%%%%%%%%%%
  % % === À DIRE À L'ORAL ===
  % % "Pourquoi ce projet ? Les outils de conception 3D classiques sont complexes, et en réalité virtuelle, 
  % % naviguer avec des manettes peut causer de la fatigue ou de la cinétose. 
  % % Nous avons donc posé deux hypothèses : d'abord, que la voix en langage naturel peut devenir la couche principale 
  % % de commande et de navigation (les mains ne servant qu'à manipuler le mobilier). 
  % % Ensuite, qu'un modèle de langage compact mais performant comme Llama 3.3 70B ou Gemini 2.5 Flash-Lite est 
  % % capable d'orchestrer cette scène 3D de manière déterministe en appelant des outils."
  % %%%%%%%%%%%%%%%%%%%%%%%%%%%%%%%%%%%%%%%%%%
\end{frame}


% ==========================================
% SLIDE 3 : Architecture Système
% ==========================================
\begin{frame}{🏗️ Architecture Générale}
  \begin{center}
    % Petit diagramme textuel
    \textbf{Client Unity (VR / PC)} $\longleftrightarrow$ \textbf{Backend FastAPI (Python)} $\longleftrightarrow$ \textbf{MongoDB}
  \end{center}

  \begin{itemize}
    \item \textbf{Unity Client} : Capture l'audio (Push-To-Talk), gère la physique, l'instanciation des prefabs et le Rig Joueur/Mains.
    \item \textbf{FastAPI Backend} : Orchestre le pipeline IA (STT $\rightarrow$ Agent LLM $\rightarrow$ TTS) et expose des endpoints de polling.
    \item \textbf{Pourquoi du Polling (1.5s - 2s) ?} 
      \begin{itemize}
        \item Le Wi-Fi des casques VR autonomes est instable.
        \item Reconnexion automatique et transparente sans couper la session utilisateur.
      \end{itemize}
  \end{itemize}

  % %%%%%%%%%%%%%%%%%%%%%%%%%%%%%%%%%%%%%%%%%%
  % % === À DIRE À L'ORAL ===
  % % "Sur le plan architectural, le système est divisé en deux blocs communicants. 
  % % D'un côté, le client Unity gère le rendu 3D, les mains rigguées, la physique et l'enregistrement audio. 
  % % De l'autre, le backend en FastAPI Python gère l'intelligence : la transcription de la parole, le raisonnement de l'agent LLM et la synthèse vocale.
  % % Plutôt que d'utiliser des WebSockets, nous avons opté pour un mécanisme de polling adaptatif toutes les 1,5 à 2 secondes. 
  % % Cela résout les instabilités courantes du Wi-Fi en VR : si la liaison coupe, le client se reconnecte silencieusement sans planter."
  % %%%%%%%%%%%%%%%%%%%%%%%%%%%%%%%%%%%%%%%%%%
\end{frame}


% ==========================================
% SLIDE 4 : L'Avatar Vocal & Rigging
% ==========================================
\begin{frame}{🧍 L'Avatar IA : Rigging et Machine d'États}
  \begin{columns}
    \begin{column}{0.5\textwidth}
      \textbf{Rigging Humanoïde (Mixamo Ch41) :}
      \begin{itemize}
        \item Squelette complet standardisé dans Unity.
        \item Suivi visuel du joueur (HUD WorldSpace attaché).
      \end{itemize}
      \vspace{0.2cm}
      \textbf{Interruption Vocale en Temps Réel :}
      \begin{itemize}
        \item Coupe la requête HTTP en cours.
        \item Stoppe la voix de l'IA instantanément.
        \item Fluidité de conversation.
      \end{itemize}
    \end{column}
    
    \begin{column}{0.5\textwidth}
      \textbf{Les 5 États d'Animation :}
      \begin{enumerate}
        \item \textbf{Sleep} : Inactif au démarrage.
        \item \textbf{WakeUp} : Transition au premier mot.
        \item \textbf{Idle} : Attente active (respiration).
        \item \textbf{Thinking} : LLM en cours (tête penchée, bras croisés).
        \item \textbf{Talking} : Edge-TTS actif (mouvements synchronisés).
      \end{enumerate}
    \end{column}
  \end{columns}

  % %%%%%%%%%%%%%%%%%%%%%%%%%%%%%%%%%%%%%%%%%%
  % % === À DIRE À L'ORAL ===
  % % "Pour donner une dimension humaine à l'assistant, l'utilisateur fait face à un avatar 3D animé. 
  % % Nous avons importé un modèle humanoïde Mixamo riggué et créé une machine d'états d'animation synchronisée avec le cycle de vie de la voix.
  % % L'avatar passe de l'état Sleep à WakeUp, puis Idle. Quand l'IA réfléchit, l'avatar croise les bras (Thinking). 
  % % Quand il parle, les lèvres bougent (Talking). 
  % % Un point technique important est le support de l'interruption : si le joueur reparle pendant que l'IA s'exprime, 
  % % la requête et l'audio se coupent immédiatement pour écouter la nouvelle instruction."
  % %%%%%%%%%%%%%%%%%%%%%%%%%%%%%%%%%%%%%%%%%%
\end{frame}


% ==========================================
% SLIDE 5 : Le Pipeline Vocal Résilient
% ==========================================
\begin{frame}{🎙️ Pipeline en Cascade à 3 Niveaux}
  \begin{table}
    \centering
    \begin{tabular}{llll}
      \toprule
      \textbf{Étape} & \textbf{Niveau 1 (Principal)} & \textbf{Niveau 2 (Secours)} & \textbf{Niveau 3 (Ultime)} \\
      \midrule
      \textbf{STT} (Parole $\rightarrow$ Texte) & Deepgram Nova-3 & Groq Whisper & Gemini Flash-Lite \\
      \textbf{LLM} (Raisonnement) & Groq Llama 3.3 70B & Groq Llama 3.1 8B & Gemini Flash-Lite \\
      \textbf{TTS} (Texte $\rightarrow$ Parole) & Edge-TTS & gTTS (local) & -- \\
      \bottomrule
    \end{tabular}
  \end{table}

  \begin{itemize}
    \item \textbf{Tolérance aux pannes} : Bascule automatique si un quota ou une API échoue.
    \item \textbf{Détection des effets de bord} : Empêche l'instanciation de doublons si un timeout survient après l'appel d'outil.
  \end{itemize}

  % %%%%%%%%%%%%%%%%%%%%%%%%%%%%%%%%%%%%%%%%%%
  % % === À DIRE À L'ORAL ===
  % % "Pour que le projet soit utilisable gratuitement et sans interruption, nous avons conçu un système de cascades d'API. 
  % % Pour chaque brique (transcription, intelligence, synthèse vocale), nous avons trois niveaux de secours. 
  % % Par exemple, pour la transcription, nous utilisons Deepgram pour sa latence de 300 millisecondes. 
  % % S'il y a une erreur réseau ou de quota, le serveur bascule automatiquement sur Groq Whisper, puis sur Gemini. 
  % % De plus, le backend intègre un système anti-doublon pour éviter qu'un meuble soit généré deux fois si l'API a mis trop de temps à répondre."
  % %%%%%%%%%%%%%%%%%%%%%%%%%%%%%%%%%%%%%%%%%%
\end{frame}


% ==========================================
% SLIDE 6 : Algorithmes de Placement
% ==========================================
\begin{frame}{📐 Algorithmes de Placement et Positionnement}
  \begin{itemize}
    \item \textbf{Détection de la position du joueur (Spatialisation)} :
      \begin{itemize}
        \item Unity envoie les coordonnées du joueur $(x, z)$ et son angle à chaque requête.
        \item Calcul trigonométrique inversé du point de spawn (3 mètres devant le joueur).
        \item Utilisation de \texttt{contextvars} côté FastAPI pour isoler les requêtes de manière thread-safe.
      \end{itemize}
    \item \textbf{Raycasting intelligent selon le type d'objet} :
      \begin{itemize}
        \item \textbf{Sol} (\textit{floor\_aware}) : Poser les chaises/tables exactement au sol en ignorant les autres colliders.
        \item \textbf{Plafond} (\textit{ceiling\_aware}) : Fixer les lampes en hauteur et orienter la lumière vers le bas.
        \item \textbf{Mur} (\textit{wall\_aware}) : Aligner les objets sur les parois verticales selon leur vecteur normal.
      \end{itemize}
    \item \textbf{Smart Offset (Évitement de collision)} :
      \begin{itemize}
        \item Test automatique de 5 positions candidates (centre, droite, gauche, devant, derrière).
        \item \texttt{Physics.OverlapSphere} détecte la présence de mobilier pour éviter le chevauchement.
      \end{itemize}
  \end{itemize}

  % %%%%%%%%%%%%%%%%%%%%%%%%%%%%%%%%%%%%%%%%%%
  % % === À DIRE À L'ORAL ===
  % % "Sous le capot, le placement des objets repose sur plusieurs algorithmes. 
  % % D'abord, le backend calcule l'endroit où faire apparaître l'objet en fonction de l'orientation et de la position de la caméra du joueur. 
  % % Ensuite, Unity utilise le Raycasting pour coller les objets aux bonnes surfaces : le sol pour le mobilier, le plafond pour les lampes avec une lumière dynamique orientée vers le bas, et les murs pour les appliques.
  % % Enfin, l'algorithme de Smart Offset effectue un test de collision par sphère. 
  % % Si l'utilisateur demande deux chaises d'affilée sans bouger, la seconde sera automatiquement décalée à côté de la première au lieu d'apparaître à l'intérieur."
  % %%%%%%%%%%%%%%%%%%%%%%%%%%%%%%%%%%%%%%%%%%
\end{frame}


% ==========================================
% SLIDE 7 : Fonctionnalités Avancées
% ==========================================
\begin{frame}{🚀 Fonctionnalités Clés et Démo}
  \begin{columns}
    \begin{column}{0.5\textwidth}
      \textbf{Téléportation Vocale Intelligente :}
      \begin{itemize}
        \item Transition douce avec fondu au noir (Fade out/in).
        \item Déplacement instantané devant la maison ou la pièce demandée.
      \end{itemize}
      \vspace{0.2cm}
      \textbf{Tableau de Bord Immersif :}
      \begin{itemize}
        \item Canvas flottant en WorldSpace UI.
        \item Affichage en direct des statistiques (latence, objets placés).
      \end{itemize}
    \end{column}
    
    \begin{column}{0.5\textwidth}
      \textbf{RAG Documentaire Réglementaire :}
      \begin{itemize}
        \item Recherche BM25 ultra-légère dans les normes DTU.
        \item Exemple : \textit{"Quelle est la taille minimale d'une chambre ?"}
      \end{itemize}
      \vspace{0.2cm}
      \textbf{Persistance MongoDB :}
      \begin{itemize}
        \item Toutes les coordonnées $(x, y, z, \text{rotation})$ et statistiques sont stockées en base.
        \item Reconstitution complète de la scène au redémarrage.
      \end{itemize}
    \end{column}
  \end{columns}

  % %%%%%%%%%%%%%%%%%%%%%%%%%%%%%%%%%%%%%%%%%%
  % % === À DIRE À L'ORAL ===
  % % "Pour compléter l'expérience, nous avons ajouté plusieurs fonctionnalités avancées. 
  % % La téléportation vocale utilise un effet de fondu pour éviter le mal des transports en VR. 
  % % Nous avons également un tableau de bord flottant en WorldSpace qui affiche les métriques de la session.
  % % Le système comprend aussi un RAG local basé sur l'algorithme BM25, qui permet à l'utilisateur de poser des questions sur les normes de construction françaises, comme les dimensions minimales d'une chambre.
  % % Et enfin, l'intégration de MongoDB permet de sauvegarder l'état complet du monde : si on ferme et relance Unity, toute la scène est fidèlement reconstruite."
  % %%%%%%%%%%%%%%%%%%%%%%%%%%%%%%%%%%%%%%%%%%
\end{frame}


% ==========================================
% SLIDE 8 : Conclusion et Perspectives
% ==========================================
\begin{frame}{🎯 Conclusion et Perspectives}
  \begin{itemize}
    \item \textbf{Livrables Réalisés} :
      \begin{itemize}
        \item Code source client (Unity) \& serveur (FastAPI) documenté.
        \item Base de données MongoDB configurée via Docker.
        \item Vidéo de démonstration en VR intégrée au README GitHub.
      \end{itemize}
    \item \textbf{Résultats} :
      \begin{itemize}
        \item Temps de réponse moyen (voix en entrée $\rightarrow$ meuble placé) inférieur à 4 secondes.
        \item Système robuste et résilient face aux coupures d'API ou de réseau.
      \end{itemize}
    \item \textbf{Perspectives} :
      \begin{itemize}
        \item Intégration de modèles 3D génératifs locaux pour créer des prefabs à la demande.
        \item Multi-joueurs collaboratif où plusieurs architectes conçoivent à la voix dans le même espace.
      \end{itemize}
  \end{itemize}

  % %%%%%%%%%%%%%%%%%%%%%%%%%%%%%%%%%%%%%%%%%%
  % % === À DIRE À L'ORAL ===
  % % "En conclusion, tous nos livrables sont opérationnels et accessibles sur GitHub, y compris le code source et la vidéo de démonstration. 
  % % Notre pipeline affiche une latence moyenne de moins de 4 secondes, ce qui est extrêmement fluide pour une interaction vocale complexe. 
  % % Pour les perspectives d'évolution, on pourrait imaginer l'intégration de modèles génératifs de mesh 3D pour créer des meubles sur mesure à la demande, ou encore l'implémentation d'un mode collaboratif multijoueur. 
  % % Je vous remercie pour votre attention et je suis disponible pour répondre à vos questions."
  % %%%%%%%%%%%%%%%%%%%%%%%%%%%%%%%%%%%%%%%%%%
\end{frame}

\end{document}
